% ===================================================================
%  TABLICA Nr 13 — Hotel. — Restauracja. — Kawiarnia.
%  Meme regime que les precedents. Tableau a UNE SEULE COLONNE, avec
%  DEUX NOTES a deux rangs differents : t13-noto-1 ouvre le fichier,
%  avant meme l'apparat, et renvoie au tableau 6 pour le detail de la
%  chambre ; t13-noto-2 se place entre le 8 et l'intertitre du
%  dejeuner, et dit au maitre qu'il peut varier le menu avec la liste
%  de plats du vocabulaire. Les deux sont en gras ET en italique dans
%  le fac-simile ; on reproduit les deux.
%
%  LE 9 EST UN MENU, ET LE GRAS Y COURT SUR LES PLATS AVEC LEUR
%  PONCTUATION : « \VUgras{butro,} », « \VUgras{legum-disho,
%  spinati} ». Les traductions sortent la virgule du gras — le
%  neerlandais le fait — parce que la regle des colonnes est le gras
%  sur le terme seul. Le fac-simile ido, lui, ne bouge pas.
%
%  ET « tripi segun la modo di Caen » RESTE UN PLAT DE CAEN. On ne
%  cherche pas d'equivalent local : le menu est un document sur ce
%  qu'on mangeait dans cet hotel-la.
% ===================================================================

%%K t13-noto-1 noto
\VUnotes{143.2mm}{%
\textit{\VUgras{(*) Co do szczegółów pokoju, zobacz tablicę Nr 6.}}%
}

%%K t13-apar-1 apar
\VUsaut{3.0mm}
\VUtitre{15.0pt}{60}{TABLICA Nr 13}
\VUsaut{1.6mm}
\VUpk{10.2pt}{40}{Hotel. --- Restauracja.}
\VUsaut{2.0mm}
\VUpk{10.2pt}{40}{Kawiarnia.}
\VUsaut{3.4mm}

%%K t13-01-1 p
1. --- Przybywszy wczoraj wieczorem do Royan, zamieszkałem w «\textit{Hotelu Oceanu}»,
który polecił mi przyjaciel.

%%K t13-01-2 p
Dano mi piękny \VUgras{pokój} \textsuperscript{(2)}, na drugim \VUgras{piętrze}
\textsuperscript{(1)}, gdzie bardzo dobrze spałem (*).

%%K t13-02-1 p
2. --- Dziś rano, wychodząc ze swojego pokoju, do którego \VUgras{pokojówka}
\textsuperscript{(3)} przyniosła śniadanie, wszedłem do \VUgras{palarni}
\textsuperscript{(5)}, która służy też za \VUgras{czytelnię} \textsuperscript{(4)}, żeby
czytać \VUgras{gazety} \textsuperscript{(6)}, paląc \VUgras{papierosa}
\textsuperscript{(8)}. Po lekturze zjechałem \VUgras{windą} \textsuperscript{(9)} i
poszedłem przejść się po mieście.

%%K t13-03-1 p
3. --- Ponieważ zapomniałem zapytać \VUgras{urzędnika} \textsuperscript{(12)} hotelu, o
której godzinie podaje się posiłki przy \VUgras{wspólnym stole} \textsuperscript{(11)},
wróciłem wcześnie i skorzystałem z tego, żeby wykąpać się w \VUgras{łazience}
\textsuperscript{(10)} hotelu. Potem poszedłem znowu do \VUgras{kasy}
\textsuperscript{(15)}, żeby zatelefonować do jednego z moich przyjaciół. Ale
\VUgras{telefon} \textsuperscript{(13)} był zajęty; widząc moje zakłopotanie,
\VUgras{kasjerka} \textsuperscript{(14)}, która zapisywała \VUgras{rachunek}
\textsuperscript{(17)} w \VUgras{księdze podróżnych} \textsuperscript{(16)}, poprosiła
\VUgras{portiera} \textsuperscript{(18)}, żeby posłał \VUgras{boya}
\textsuperscript{(19)} załatwić moje zlecenie \VUgras{na rowerze}
\textsuperscript{(20)}.

%%K t13-04-1 p
4. --- Podziękowałem jej i wyszedłem, żeby dać napiwek \VUgras{tragarzowi}
\textsuperscript{(24)} (robotnikowi od dźwigania), który przywiózł z dworca na swoim
\VUgras{ręcznym wózku} \textsuperscript{(25)} moje \VUgras{pudło na kapelusz}
\textsuperscript{(26)}, \VUgras{moją walizkę} \textsuperscript{(27)} i \VUgras{moją
torbę podróżną} \textsuperscript{(28)}. \VUgras{Listonosz} \textsuperscript{(21)},
który właśnie przybył, dał mi \VUgras{list} \textsuperscript{(22)}.

%%K t13-05-1 p
5. --- Zostałem jeszcze przez jakiś czas na progu drzwi, blisko \VUgras{skrzynki na
listy} \textsuperscript{(23)}, żeby popatrzeć na ruch na ulicy. \VUgras{Konduktor}
\textsuperscript{(30)} \VUgras{omnibusu} \textsuperscript{(29)} ładował na swój wóz
\VUgras{kufer} \textsuperscript{(31)} \VUgras{podróżnego} \textsuperscript{(32)}, z
którym żegnał się \VUgras{właściciel hotelu} \textsuperscript{(33)}. Młody
\VUgras{roznosiciel telegramów} \textsuperscript{(35)} bez pośpiechu przyniósł
\VUgras{telegram} \textsuperscript{(34)} dla klienta hotelu; \VUgras{turysta}
\textsuperscript{(36)} ze swoim \VUgras{aparatem fotograficznym} \textsuperscript{(38)}
przez ramię sprawdzał swój \VUgras{przewodnik} \textsuperscript{(37)}.

%%K t13-tit-1 sub
\VUsaut{4.0mm}
\VUpk{10.2pt}{40}{Pora obiadu.}
\VUsaut{3.0mm}

%%K t13-06-1 p
6. --- Przed kratą beztroski \VUgras{posłaniec} \textsuperscript{(39)} drzemał między
swoim \VUgras{hakiem do noszenia} \textsuperscript{(40)} a swoją \VUgras{skrzynką do
czyszczenia butów} \textsuperscript{(41)}, podczas gdy młoda \VUgras{praczka}
\textsuperscript{(42)}, która niosła \VUgras{kosz} \textsuperscript{(43)} bielizny,
śmiała się z uroczystej miny \VUgras{kierownika sali} \textsuperscript{(44)}, który stał
przy drzwiach.

%%K t13-07-1 p
7. --- Widząc \VUgras{kucharza} \textsuperscript{(45)}, który wyszedł ze swojej
\VUgras{kuchni} \textsuperscript{(47)} położonej w \VUgras{suterenie}
\textsuperscript{(48)}, żeby posłać \VUgras{pomywacza} \textsuperscript{(46)} załatwić
zlecenie, przypomniałem sobie, że godzina obiadu jest bliska, i wszedłem do
restauracji, przechodząc przez podwórze, na którym \VUgras{automobilista}
\textsuperscript{(50)} stawiał swój \VUgras{samochód} \textsuperscript{(49)}, podczas
gdy \VUgras{mechanik} \textsuperscript{(52)} naprawiał, według wskazówek
\VUgras{rowerzysty} \textsuperscript{(53)}, \VUgras{trójkołowiec naftowy}
\textsuperscript{(51)}, który właśnie uległ wypadkowi.

%%K t13-08-1 p
8. --- Zapytałem młodego \VUgras{boya} \textsuperscript{(54)}, który niósł
\VUgras{szklanki piwa} \textsuperscript{(55)}, czy wspólny stół jest pod werandą, a
wobec jego twierdzącej odpowiedzi zająłem miejsce przy jednym ze stołów i kazałem sobie
dać wyborny posiłek.

%%K t13-noto-2 noto
\VUnotes{143.2mm}{%
\textit{\VUgras{(*) Za pomocą listy dań wprowadzonej do słownika nauczyciel będzie mógł
łatwo urozmaicać poniższe menu.}}%
}

%%K t13-tit-2 sub
\VUsaut{4.0mm}
\VUpk{10.2pt}{40}{Wyborny obiad.}
\VUsaut{3.0mm}

%%K t13-09-1 p
9. --- Kelner przyniósł mi menu (*) obiadu i listę win. Wybrałem i zamówiłem jako
\VUgras{przystawkę krewetki} z \VUgras{masłem}, a jako danie wstępne \VUgras{flaki na
sposób z Caen}. Nie jadłem \VUgras{ryby}, ale poprosiłem o porcję baraniego
\VUgras{udźca}, a jako \VUgras{danie jarzynowe, szpinak} ze śmietaną. Potem, ponieważ
żadna z \VUgras{przekąsek} mi się nie podobała, kazałem przynieść rozmaite desery,
\VUgras{ser}, \VUgras{owoce} i \VUgras{biszkopty}. Dodałem do tego jeszcze wyborne
\VUgras{wino} Saint-Émilion, i bardzo zadowolony ze swojego posiłku opuściłem stół,
dawszy \VUgras{kelnerowi} \textsuperscript{(57)} piękny \VUgras{napiwek}
\textsuperscript{(59)}.

%%K t13-10-1 p
10. --- Widząc mnie gotowego do odejścia, \VUgras{zarządca} \textsuperscript{(60)}
zaproponował mi, żebym udał się na balkon podziwiać wspaniałą panoramę \VUgras{Oceanu}
\textsuperscript{(62)}. W dali zaczynało się widzieć rybackie \VUgras{łódeczki}
\textsuperscript{(63)}, \VUgras{żaglowce} \textsuperscript{(64)} i ogromny
\VUgras{parowiec pocztowy} \textsuperscript{(65)}, którego czarna sylwetka odcina się na
białych \VUgras{chmurach} \textsuperscript{(67)} \VUgras{horyzontu}
\textsuperscript{(66)}. Lekki wiaterek kazał falować \VUgras{sztandarowi}
\textsuperscript{(70)} zatkniętemu nad \VUgras{szyldem} \textsuperscript{(69)}
\VUgras{hali} \textsuperscript{(68)}.

%%K t13-tit-3 sub
\VUsaut{3.0mm}
\VUpk{10.2pt}{40}{Rozrywki.}
\VUsaut{3.0mm}

%%K t13-11-1 p
11. --- Widowisko było tak ciekawe, że postanowiłem wejść jutro na taras kawiarni,
zabierając ze sobą \VUgras{lornetkę teatralną} \textsuperscript{(71)} albo
\VUgras{lunetę} \textsuperscript{(72)}.

%%K t13-12-1 p
12. --- Opuszczając balkon, udałem się do kawiarni hotelu, żeby wypić \VUgras{napój}
\textsuperscript{(73)}, grając \VUgras{partię bilardu} \textsuperscript{(75)}. Bez
zatrzymania przeszedłem przez salę kawiarnianą, w której wielu konsumentów grało w
\VUgras{szachy} \textsuperscript{(78)}, w \VUgras{warcaby} \textsuperscript{(79)}, w
\VUgras{kostki domina} \textsuperscript{(80)}, w \VUgras{tryktraka}
\textsuperscript{(81)} albo w \VUgras{karty} \textsuperscript{(82)}, i wprost wszedłem
do \VUgras{sali bilardowej} \textsuperscript{(74)}.

%%K t13-13-1 p
13. --- Kiedy partia się skończyła, zszedłem na parter i poprosiłem kelnera o
\VUgras{kopertę} \textsuperscript{(84)}, \VUgras{papier} \textsuperscript{(83)},
\VUgras{znaczek pocztowy} \textsuperscript{(85)}, \VUgras{pióro} \textsuperscript{(87)} i
\VUgras{atrament} \textsuperscript{(86)}, żeby napisać list.

%%K t13-13-2 p
Potem przywołałem \VUgras{woźnicę} \textsuperscript{(92)}, którego \VUgras{powóz}
\textsuperscript{(88)} zatrzymał się przed kawiarnią. Zapytałem go o cenę według godzin
albo za kurs, a kiedy zgodziliśmy się co do ceny, otworzył mi \VUgras{drzwiczki powozu}
\textsuperscript{(89)} i usadowił mnie. Wrócił na \VUgras{kozioł} \textsuperscript{(91)},
szybko ruszył konie \VUgras{bacikiem} \textsuperscript{(93)} i wyjechaliśmy na uroczą
przejażdżkę.
\VUsaut{6.0mm}
\VUfilet{22mm}
